{{دیوان
| عنوان = جاوید کے نام (دیارِ عشق میں اپنا مقام پیدا کر)
| ماخذ = ویکی ماخذ
| ماخذ_لنک = https://ur.wikisource.org/wiki/%D8%AC%D8%A7%D9%88%DB%8C%D8%AF_%DA%A9%DB%92_%D9%86%D8%A7%D9%85_%28%D8%AF%DB%8C%D8%A7%D8%B1%D9%90_%D8%B9%D8%B4%D9%82_%D9%85%DB%8C%DA%BA_%D8%A7%D9%BE%D9%86%D8%A7_%D9%85%D9%82%D8%A7%D9%85_%D9%BE%DB%8C%D8%AF%D8%A7_%DA%A9%D8%B1%29
| مقابلہ = اقبال سائبر لائبریری (اقبال اکادمی پاکستان)
| مقابلہ_لنک = https://www.iqbalcyberlibrary.net
}}
=== (لندن میں اُس کے ہاتھ کا لِکھّا ہُوا پہلا خط آنے پر) ===

<poem>
دیارِ عشق میں اپنا مقام پیدا کر
نیا زمانہ، نئے صبح و شام پیدا کر

خدا اگر دل فطرت شناس دے تجھ کو
سکُوتِ لالہ و گُل سے کلام پیدا کر

اُٹھا نہ شیشہ گرانِ فرنگ کے احساں
سفالِ ہند سے مِینا و جام پیدا کر

میں شاخِ تاک ہوں، میری غزل ہے میرا ثمر
مرے ثمر سے میء لالہ فام پیدا کر

مرا طریق امیری نہیں، فقیری ہے
خودی نہ بیچ، غریبی میں نام پیدا کر!
</poem>
